% bread-types-ar.tex — Arabic edition of ch. 8, bread types
%
% Original: The Sourdough Framework by Hendrik Kleinwächter,
% the-sourdough-framework.com, CC BY-SA 4.0. The Arabic edition inherits
% CC BY-SA 4.0.
%
% Phase 28R R2 unit 10 reconstruction. Body text recovered verbatim from
% the R0 recovery reference (recovery-r0/text/08-bread-types.txt,
% mined from the Phase 23 website build of the accepted Arabic sources)
% with the accepted edition book-ar-fixed.pdf (b78c2a28b574288801d8000a7249af3a)
% as the fidelity authority; no text retranslated from the English.
% Float and short captions recovered from the accepted edition's
% renders (pp. 61–71) and its list-of-floats pages (short captions
% قائمة الجداول / قائمة المخططات الانسيابية / قائمة الأشكال).
%
% Direction treatments verified against the accepted renders:
% tables 8.1/8.2 are mirrored with TRT cells and PLAIN \thead header
% cells (probe10d/probe10e, 2026-11-15: the accepted headers' centered
% appearance is an artifact of RTL content in l columns — the rows are
% the exact upstream mirrors; tables-ar/ header comments carry the
% geometry); the recipe tabular
% keeps the upstream {r@{}rl@{}} spec verbatim (NOT mirrored); the
% flatbread flowchart is NOT mirrored (اخلط المكونات leftmost →
% انتظر حتى يجهز العجين middle → اخبز الخبز على الموقد rightmost).
% Body ranges are literal-text LTR islands per the unit 9 dash-font
% convention (\LR{\textenglish{6 h\textendash 12 h}} — a bare
% \textendash inside \LR{} picks Amiri, probe10b N; the accepted
% renders every body dash in Pagella). Cross-chapter refs are kept as
% literal text (القسم~\LR{5.1} ``الحساب الخبازي'' for upstream
% \ref/\nameref{sec:bakers-math}; فصل الخبز الحامض القمحي for upstream
% \nameref{ch:wheat-sourdough}) since those labels are not yet built in
% unit builds — the undefined-refs pending set on the R3 agenda.
% Documented residual (paren-font): ASCII parens in Arabic-font context
% render Pagella here vs Amiri in the accepted edition (lua fontsplit
% remap, unfixable at the content layer; R3 harmonization agenda).

\chapter{أنواع الخبز}\label{ch:bread-types}

\begin{quoting}
ستتعلم في هذا الفصل عن أنواع الخبز المختلفة ومزايا كلٍّ منها وعيوبه. ستجد هنا أيضًا وصفات بسيطة جدًا للخبز المُسطَّح ولخبز القالب. فالأول هو على الأرجح أيسر أنواع الخبز وأقلها جهدًا يمكن صنعه، أما الثاني فيتطلب جهدًا يزيد قليلًا. وللخبز القائم بذاته فصل خاص به نظرًا لتعقيده الأكبر.
\end{quoting}

\section{المقدمة}\label{sec:intro}

في هذا القسم نصنّف الخبز بحسب تقنيات خبزه. سيختلف المظهر والطعم بالطبع، لكن يمكنك أن تحصل على خبز ممتاز بكل تقنية منها. بعض أنواع الخبز تتطلب استثمارًا وتقنية، كما يعرض الجدول~\LR{\ref{tab:bread-types-comparison}}. والخبز المُسطَّح هو على الأرجح أيسر أنواع الخبز وأقلها جهدًا يمكن صنعه. فإن كنت شخصًا مشغولًا و/أو لا تملك فرنًا، فقد يكون هذا بالضبط نوع الخبز الذي ينبغي أن تفكر فيه.

\begin{table}[!htb]
\centering
% table-overview-bread-types-ar.tex — Arabic edition of upstream
% tables/table-overview-bread-types.tex (ch. 8, جدول 8.1)
%
% Original: The Sourdough Framework by Hendrik Kleinwächter,
% the-sourdough-framework.com, CC BY-SA 4.0. The Arabic edition inherits
% CC BY-SA 4.0.
%
% RTL mirror per the settled table mechanics: the column order is
% reversed in the declared LTR tabular (upstream {@{}llll@{}} keeps
% {@{}llll@{}} under reversal) so طريقة الخبز renders rightmost and
% القائم بذاته leftmost, matching the accepted edition's render of
% p. 61: data columns القائم بذاته (x154.5–197.5 header), قالب الخبز
% (222.8–266.3), المُسطَّح (317.2–347.4), label column rightmost (rule
% starts 144.7, ends 454.4).
%
% Header geometry (measured against accepted p. 61): the span header
% \multicolumn{3}{c}{\textbf{نوع الخبز}} centers at 266.7 — the exact
% center of the three data columns INCLUDING the trailing \tabcolsep
% (144.7–388.8), confirming the {@{}llll@{}} spec; its
% \cmidrule(lll){1-3} spans 150.6–388.8 (left trim only, 5.9 pt). The
% three subheaders are PLAIN \thead cells in the l columns — the exact
% upstream mirror (upstream: `& \textbf{Flat} & \textbf{Loaf pan} &
% \textbf{Free standing}`). PROBE10D FINDING (2026-11-15): the all-
% plain row reproduces the accepted ink positions EXACTLY (col1
% 154.5–197.5, col2 222.8–266.3, col3 317.2–347.4, i.e. +9.8/+3.6/
% +35.4 from the column edges); their centered appearance in the
% accepted render is an artifact of RTL content in l columns, not
% \multicolumn centering. A \multicolumn{1}{c} wrapper on col1 renders
% its ink +2.9 pt right of the accepted position (probe10c T1/T2/T3:
% \thead vs \bfseries and @{}c@{} padding variants all identical);
% on col2/col3 the {c} wrapper happens to render identically to plain
% (probe10d T6), but the plain form is kept as the faithful upstream
% mirror. The label column carries regular Amiri weight (no \thead)
% per the accepted render.
%
% SPAN HEADER (PROBE10F FINDING, 2026-11-15): the span cell is
% \thead{نوع الخبز}, NOT \textbf — the accepted toprule→cmidrule
% distance is 28.4 pt with the span text 6.7 pt below the toprule;
% \textbf gives 22.5 pt / 3.8 pt (the nested-tabular row of \thead
% supplies the extra height). \thead leaves the span text's position
% relative to the table unchanged (probe10f W1/W2: text center minus
% table center identical), so the x-extents stay exact.
%
% CELL DIRECTION (BINDING, unit 9 verdict): every non-empty cell
% wrapped in {\textdir TRT ...} at the content layer — the sty's
% \insert@column \textdir TRT patch is a silent no-op. Digit runs ride
% per-token \LR islands (probe10 H form: {\textdir TRT \LR{60} دقيقة}).
% Documented residual: the accepted edition renders digits in Pagella
% with the Arabic in Amiri; the delivered cells match (island digits
% ride the lua remap to Pagella).

\begin{tabular}{@{}llll@{}}
\toprule
\multicolumn{3}{c}{\thead{{\textdir TRT نوع الخبز}}} & \\
\cmidrule(lll){1-3}
{\thead{{\textdir TRT القائم بذاته}}} & {\thead{{\textdir TRT قالب الخبز}}} & {\thead{{\textdir TRT المُسطَّح}}} & \\
\midrule
{\textdir TRT الفرن} & {\textdir TRT الفرن} & {\textdir TRT المقلاة أو النار أو الشوّاية} & {\textdir TRT طريقة الخبز} \\
{\textdir TRT \LR{60} دقيقة} & {\textdir TRT \LR{5} دقائق} & {\textdir TRT \LR{3} دقائق} & {\textdir TRT زمن العمل} \\
{\textdir TRT الدقيق الغلوتيني} & {\textdir TRT جميع الأنواع} & {\textdir TRT جميع الأنواع} & {\textdir TRT أنواع الدقيق} \\
{\textdir TRT صعبة} & {\textdir TRT سهلة} & {\textdir TRT سهلة جدًا} & {\textdir TRT درجة الصعوبة} \\
{\textdir TRT مرتفعة} & {\textdir TRT متوسطة} & {\textdir TRT منخفضة} & {\textdir TRT الكلفة} \\
\bottomrule
\end{tabular}

\caption[أنواع الخبز المختلفة]{نظرة عامة على أنواع الخبز المختلفة ودرجة تعقيد كلٍّ منها.}\label{tab:bread-types-comparison}
\end{table}

\section{الخبز المُسطَّح}\label{sec:flatbread}

الخبز المُسطَّح هو على الأرجح أبسط خبز حامض يمكن صنعه. وصنعه لا يتطلب فرنًا؛ كل ما تحتاجه هو موقد.

\begin{figure}[!htb]
\centering
\includegraphics[width=\textwidth]{flat-breads-selection}
\caption[تشكيلة من الخبز المُسطَّح]{تشكيلة منتقاة من الخبز المُسطَّح المختلف المصنوع من الخبز الحامض. من اليسار إلى اليمين: تورتيلا قمحية \textenglish{(tortilla)}، جاودار، سبلت، وذرة.}
\end{figure}

هذا النوع من الخبز بسيط للغاية في صنعه إذ يمكنك تجاوز كثير من التقنية التي يتطلبها عادةً صنع عجائن القمح. ويمكن صنع الخبز المُسطَّح بكل أنواع الدقيق. بل يمكنك حتى استخدام دقيق خالٍ من الغلوتين، مثل دقيق الذرة أو الأرز، لصنع العجين. ولكي تجعل الخبز المُسطَّح أكثر انتفاخًا قليلًا، يمكنك استخدام قليل من دقيق القمح. فالغلوتين الذي يتكوّن سيحبس الغازات. وخلال الخَبز ستنفخ هذه الغازات العجين.

وحيلة أخرى لتحسين ملمس الخبز المُسطَّح هي صنع عجين سائل جدًا. فكثير من الماء سيتبخر أثناء عملية الخَبز فيصير الخبز أنفس وأكثر انتفاخًا. وإذا كان محتوى الماء لديك مرتفعًا جدًا فسيُنتج قوامًا شبيهًا بالقرص المسطّح، كما ترى في الجدول~\LR{\ref{tab:flat-bread-ingredients}}

\begin{table}[!htb]
\centering
% table-flat-bread-pancake-recipe-ar.tex — Arabic edition of upstream
% tables/table-flat-bread-pancake-recipe.tex (ch. 8, جدول 8.2)
%
% Original: The Sourdough Framework by Hendrik Kleinwächter,
% the-sourdough-framework.com, CC BY-SA 4.0. The Arabic edition inherits
% CC BY-SA 4.0.
%
% RTL mirror per the settled table mechanics: the column order is
% reversed in the declared LTR tabular (upstream {@{}lll@{}} keeps
% {@{}lll@{}} under reversal) so the label column renders rightmost
% (x414.1), الخبز المُسطَّح the rightmost data column (header
% 315.3–368.6) and الأقراص المسطّحة the leftmost (header 169.6–245.0),
% matching the accepted edition's render of p. 63 (rules 144.8–454.3 at
% y70.6/105.0/190.7).
%
% Header row (PROBE10E CORRECTION, 2026-11-15): the prior-session row
% put الأقراص المسطّحة in the LABEL column (col3) with col1 empty —
% structurally wrong. The accepted p. 63 render shows col1 =
% الأقراص المسطّحة (ink 169.6–245.0, over the pancake amounts column),
% col2 = الخبز المُسطَّح (ink 315.3–368.6) and the label column EMPTY:
% the exact mirror of the upstream header `& \textbf{Flat breads} &
% \textbf{Pancakes}`. The row is written all-plain (no \multicolumn):
% probe10e V2/V3 render identically, and probe10d showed {c} wrappers
% on first-column headers shift the ink +2.9 pt (table 8.1). The
% centered appearance is an artifact of RTL content in l columns.
%
% Header geometry (measured against accepted p. 63): الأقراص المسطّحة
% +24.8 to +100.2 from the col1 edge, الخبز المُسطَّح +33.5 to +86.8
% from the col2 edge; probe10e V3 reproduces the pancake header
% exactly.
%
% Cell construction (probe10b, verified at char level against accepted
% p. 63/p. 68):
%  - حتى 100 g (100 %) per form I: Arabic + ONE \LR island holding the
%    whole number-paren-percent run, so the island's internal LTR order
%    survives the TRT cell (bare parens BETWEEN islands flip island
%    visual order in a TRT cell — probe10 A/B/F).
%  - range cells per form K: \LR{\textenglish{5 g\textendash 20 g}
%    {\arabicfont(}\textenglish{5 \%\textendash 20 \%}}{\arabicfont)}
%    — literal-text ranges per the unit 9 dash-font convention (a bare
%    \textendash inside \LR{} picks Amiri; \textenglish wrapping yields
%    the Pagella dash the accepted edition shows).
%  - حين يزداد حجم العجين 50 % per form M: mixed Arabic + \LR{50 \%}
%    island.
% Documented residual (paren-font): the accepted edition renders the
% percent parens in Amiri (p63/p68); the lua fontsplit remap re-fonts
% ASCII punctuation to Pagella unconditionally, so our build renders
% them Pagella. Unfixable at the content layer; logged for the R3
% harmonization agenda.

\begin{tabular}{@{}lll@{}}
\toprule
{\thead{{\textdir TRT الأقراص المسطّحة}}} & {\thead{{\textdir TRT الخبز المُسطَّح}}} & \\
\midrule
{\textdir TRT \LR{\qty{100}{\gram}}} & {\textdir TRT \LR{\qty{100}{\gram}}} & {\textdir TRT الدقيق} \\
{\textdir TRT \LR{\qty{300}{\gram}} {\arabicfont(}\qty{300}{\percent}{\arabicfont)}} & {\textdir TRT حتى \LR{100 g {\arabicfont(}100 \%{\arabicfont)}}} & {\textdir TRT الماء} \\
{\textdir TRT \LR{\textenglish{5 g\textendash 20 g} {\arabicfont(}\textenglish{5 \%\textendash 20 \%}}{\arabicfont)}} & {\textdir TRT \LR{\textenglish{5 g\textendash 20 g} {\arabicfont(}\textenglish{5 \%\textendash 20 \%}}{\arabicfont)}} & {\textdir TRT البادئة} \\
{\textdir TRT \LR{\qty{2}{\gram}} {\arabicfont(}\qty{2}{\percent}{\arabicfont)}} & {\textdir TRT \LR{\qty{2}{\gram}} {\arabicfont(}\qty{2}{\percent}{\arabicfont)}} & {\textdir TRT الملح} \\
{\textdir TRT حين تظهر فقاعات على السطح} & {\textdir TRT حين يزداد حجم العجين \LR{50 \%}} & {\textdir TRT متى تخبز؟} \\
\bottomrule
\end{tabular}

\caption[وصفة الخبز المُسطَّح]{وصفة الخبز المُسطَّح أو القرص المسطّح لشخص واحد. اضاعف المقادير لتكبير حجم الحصة. راجع القسم~\LR{5.1} ``الحساب الخبازي'' لتتعلم كيف تفهم النسب المئوية وتستخدمها على نحو سليم.}\label{tab:flat-bread-ingredients}
\end{table}

للحصول على وصفة كاملة تشمل عملية صنع مثل هذا الخبز المُسطَّح، راجع القسم الفرعي~\LR{\ref{subsec:flat-bread-recipe}}

\subsection{الإطار العام للخبز المُسطَّح}\label{subsec:flat-bread-framework}

كما شرحنا أعلاه، فإن كنت في بداية طريقك فصنع الخبز المُسطَّح هو أيسر طريقة لبدء صناعة خبز رائع في المنزل. وبخطوات قليلة يمكنك التوقف عن شراء الخبز إلى الأبد. وهذا يعمل مع أي دقيق، بما في ذلك الخيارات الخالية من الغلوتين.

\begin{flowchart}[!htb]
\centering
% fig-process-flat-bread-ar.tex — Arabic edition of the flatbread
% process flowchart (ch. 8, المخطط الانسيابي 8.1)
%
% Original: The Sourdough Framework by Hendrik Kleinwächter,
% the-sourdough-framework.com, CC BY-SA 4.0. The Arabic edition inherits
% CC BY-SA 4.0.
%
% Phase 28R R2 unit 10 reconstruction. Node/edge geometry verbatim from
% the upstream source figures/fig-process-flat-bread.tex (node styles
% from figures/flowcharts_tikz.tex, which sourdough-ar.sty inputs);
% label text recovered from the accepted edition's render of p. 63
% (book-ar-fixed.pdf, b78c2a28b574288801d8000a7249af3a) — label
% recovery against a visual reference, not retranslation.
%
% Layout is NOT mirrored, per the accepted render: اخلط المكونات
% leftmost → انتظر حتى يجهز العجين middle → اخبز الخبز على الموقد
% rightmost (same reading order as the upstream English flowchart).
% Manual line splits follow the accepted render exactly: 2 lines /
% 3 lines (انتظر / حتى يجهز / العجين) / 2 lines. Upstream picture
% options (node distance=3cm, auto, every node inner sep=10, outer
% sep=0) and per-node node distance=5cm are kept verbatim; the accepted
% render shows 142.2 pt center-to-center spacing, matching 5 cm at the
% scaled node text.
%
% Arabic label text is wrapped in \textarabic{} per the established
% flowchart convention (fig-starter-process-ar.tex). Documented
% residual (same family as the earlier flowchart units): node text
% rides the scaled Amiri (10.79 pt) where the accepted edition's
% standalone-figure builds show 10.04 pt.

\begin{tikzpicture}[node distance=3cm, auto, every node/.style={inner sep=10, outer sep=0}]
  \node [start] (init) {\textarabic{اخلط}\\ \textarabic{المكونات}};
  \node [block, right of=init, node distance=5cm] (wait) {\textarabic{انتظر}\\ \textarabic{حتى يجهز}\\ \textarabic{العجين}};
  \node [success, right of=wait, node distance=5cm] (bake_bread) {\textarabic{اخبز الخبز}\\ \textarabic{على الموقد}};

  \draw [line] (init) -- (wait);
  \draw [line] (wait) -- (bake_bread);
\end{tikzpicture}

\caption[عملية صنع الخبز الحامض المُسطَّح]{عملية صنع الخبز المُسطَّح بسيطة جدًا، وتتطلب جهدًا يسيرًا. وهذا النوع من الخبز مفيد بوجه خاص للخبازين المشغولين.}\label{fig:flat-bread-process}
\end{flowchart}

هذه هي الوصفة التي أعتمدها لصنع الخبز كلما ضاق بي الوقت أو كنت في الخارج. ويمكنك الاختيار بين خيارين:
%
\begin{enumerate}
    \item خبز مُسطَّح يشبه الروتي \textenglish{(roti)} أو خبز النان \textenglish{(naan)}
    \item أقراص مسطّحة من الخبز الحامض.
\end{enumerate}

لتبدأ، جهّز بادئة الخبز الحامض لديك. فإن لم تُستخدم منذ مدة طويلة جدًا، ففكر في منحها تغذية أخرى. ولعمل ذلك، خذ ببساطة \LR{\qty{1}{\gram}} من بادئة الخبز الحامض الموجودة لديك وغذِّها بـ\LR{\qty{5}{\gram}} من الدقيق و\LR{\qty{5}{\gram}} من الماء. فإن فعلت ذلك صباحًا، ستكون بادئة الخبز الحامض جاهزة بحلول المساء. و كلما ارتفعت الحرارة أجهزت أسرع، ففكر في استخدام ماء دافئ إن كان الجو شديد البرودة حيث تسكن.

\begin{figure}[htb!]
\centering
\includegraphics[width=1.0\textwidth]{flat-bread-wheat}
\caption[خبز مُسطَّح من القمح]{خبز مُسطَّح مصنوع من دقيق القمح وحده. العجين أكثر جفافًا عند ترطيب يقارب \LR{\qty{60}{\percent}}. والعجين الأجف أصعب قليلًا في العجن. ولما كان القمح يحوي غلوتينًا أكثر، ينتفخ العجين أثناء عملية الخَبز.}
\end{figure}

بهذه الطريقة سيكون لديك نحو \LR{\qty{11}{\gram}} من بادئة الخبز الحامض جاهزة في المساء. وستحصل على الكمية المثالية لصنع عجين لشخص واحد. وإن أردت صنع خبز أكثر، فاضاعف ببساطة الكميات المبينة في الجدول~\LR{\ref{tab:flat-bread-ingredients}}.

ثم في المساء اخلط المقادير كما يعرضها الجدول. عجينك سيكون جاهزًا في الصباح. وهو عادةً جاهز بعد \LR{\textenglish{6 h\textendash 12 h}}. فإن استخدمت بادئة خبز حامض أكثر سيجهز أسرع، وبالمقابل سيستغرق وقتًا أطول إن استخدمت أقل. حاول أن تستهدف زمن تخمير يتراوح بين \LR{\textenglish{8 h\textendash 12 h}} إذ باستخدام عجينك قبل أوانه قد لا تكون النكهة بأفضل جودة. وباستخدام عجينك لاحقًا قد يصير أكثر حامضية قليلًا. والخيار الأفضل هو التجريب ورؤية ما يعجبك شخصيًا أكثر.

بعد خلط المقادير معًا غطِّ الوعاء؛ فهذا يمنع جفاف العجين ويضمن ألا تصل إليه ذباب الفاكهة. وسيساعدك الوعاء الشفاف عند البدء. يمكنك ملاحظة العجين بسهولة أكبر ورؤية متى يجهز.

\begin{figure}[htb]
\centering
\includegraphics[width=1.0\textwidth]{ethiopian-woman-checking-bread}
\caption[امرأة إثيوبية تخبز الإنجيرا]{امرأة إثيوبية تخبز إنجيرا \textenglish{(injera)} مصنوعة من دقيق التيف \textenglish{(teff)}. وفّرت الصورة \textenglish{Charliefleurene} عبر ويكيبيديا.}
\end{figure}

إن استخدمت خيار الخبز المُسطَّح بماء أقل، فانظر إلى زيادة الحجم في عجينك. ينبغي أن يكون قد زاد \LR{\qty{50}{\percent}} على الأقل. وانظر أيضًا إلى الفقاعات على جوانب وعائك.

وعند استخدام وصفة القرص المسطّح، فقِر الفقاعات على سطح عجينك. وفي الحالتين استخدم أنفك لتفحص رائحة عجينك. وبحسب النظام البيئي الدقيق لبادئة الخبز الحامض لديك سيحمل عجينك نفحات لبنية أو فاكهية أو كحولية، أو نفحات خليةً أسيتية. والاعتماد على رائحة عجينك هو أفضل وسيلة للحكم على جهوز عجينك من عدمه. فالأزمنة ليست مؤتمنة لأنها تعتمد على بادئتك وعلى الحرارة. فإن جهز عجينك قبل أوانه، يمكنك الآن نقله مباشرة إلى الثلاجة وخبزه في وقت لاحق أنسب لك. فالحرارة المنخفضة ستوقف عملية التخمير\footnote{هناك بعض الاستثناءات. ففي حالات نادرة قد تعمل بادئتك أيضًا عند درجات حرارة أدنى. فقد تكون ربّيت كائنات دقيقة تعمل بأفضل حال عند درجات الحرارة المنخفضة. ومع ذلك، التخمير يظل أبطأ كلما انخفضت الحرارة. والثلاجة تساعد فعلًا في حفظ حالة عجينك.} وسيدوم عجينك عدة أيام. وكلما طال الانتظار، صار الخبز أكثر حامضية. والثلاجة خيار رائع إن أردت أن تأخذ العجين معك في زيارة للأصدقاء. وسيعشقك الناس لأجل الأقراص المسطّحة أو الخبز المُسطَّح الطازج. وإن جاز لك الأمر، يمكنك أيضًا أن تتذوق قليلًا من عجينك الخام غير المطهو. ومن المرجح أن يكون طعمه حامضيًا نسبيًا. وأنا أفعل ذلك كثيرًا لأقيّم على نحو أفضل حالة عجائني.

\begin{figure}[!htb]
\centering
\includegraphics[width=1.0\textwidth]{injera-pancake-texture}
\caption[قرص مسطَّح من دقيق التيف]{قرص مسطّح من الخبز الحامض مصنوع من دقيق التيف. والجيوب ناتجة عن الماء المتبخر و\LR{\ch{CO2}} الذي تصنعه الكائنات الدقيقة. وفّرت الصورة \textenglish{Łukasz Nowak} عبر ويكيبيديا.}
\end{figure}

وإن كنت تشعر بالكسل أو لم يكن لديك وقت، فيمكنك أيضًا استخدام بادئة خبز حامض أقدم لصنع العجين مباشرة دون أي تغذيات أولية للبادئة. وسيعيد بادئتك النموّ داخل عجينك. تذكّر أن الخبز النهائي قد يكون أكثر حامضية قليلًا إذ قد يكون توازن الخمائر إلى البكتيريا مضطربًا. وفي الجدول~\LR{\ref{tab:flat-bread-ingredients}} أنصح باستخدام نحو \LR{\textenglish{5 \%\textendash 20 \%}} من بادئة الخبز الحامض بحسابها على الدقيق لصنع العجين. وإن أردت اتباع هذا النهج، فاستخدم نحو \LR{\qty{1}{\percent}} فقط واصنع العجين مباشرة. ومن المرجح أن يكون العجين جاهزًا بعد \num{24} ساعة، بحسب الحرارة.

وإن أردت صنع أقراص مسطّحة حلوة، فأضف إلى عجينك الآن بعض السكر وبيضًا اختياريًا. والكمية الجيدة من البيض نحو بيضة واحدة لكل \LR{\qty{100}{\gram}} من الدقيق. اقلب عجينك قليلًا وسيكون جاهزًا للاستخدام. وسيكون لديك أقراص مسطّحة حلوة لذيذة، المزيج المثالي. وب إضافة السكر الآن، تضمن أن الكائنات الدقيقة لا تملك وقتًا كافيًا لتخميره تخميرًا تامًا. ولو كنت قد أضفت السكر قبل ذلك، لما بقي أي طعم حلو بعد \num{12} ساعة.

لخبز عجينك سخّن موقدك إلى حرارة متوسطة. أضف قليلًا من الزيت إلى المقلاة. فهذا يساعد على توزيع الحرارة ويضمن طهيًا متساويًا. وبسباتولا \textenglish{(spatula)} أو بملعقة ضع عجينك في المقلاة. فإن كان عجينك في الثلاجة، فاخبزه مباشرة. لا حاجة إلى انتظار عجينك حتى يعود إلى حرارة الغرفة. وإن كان لديك غطاء، ضعيه على مقلاتك. فالغطاء يساعد على طهي عجينك من الأعلى. وسيتبخر الماء ويجول ويسخّن سطح العجين. عند صنع الخبز المُسطَّح، اجعل العجين بسماكة نحو \LR{\qty{1}{\cm}}. وعند استخدام خيار القرص المسطّح، اختر نحو \LR{\textenglish{0.1 cm\textendash 0.5 cm}} بحسب ما يعجبك.

\begin{figure}[htb]
\centering
\includegraphics[width=1.0\textwidth]{einkorn-crumb}
\caption[لُبّ الخبز من القمح الأحادي]{لُبّ خبز مُسطَّح مصنوع من دقيق القمح الأحادي. فالقمح الأحادي خالٍ من الغلوتين تقريبًا، ولذلك لا يحبس قدر \LR{\ch{CO2}} الذي يحبسه عجين قمحي. ولكي تجعل العجين أنفس استخدم ماء أكثر أو فكر في إضافة قمح أكثر إلى خليط عجينك.}
\end{figure}

بعد \LR{\textenglish{2 min\textendash 4 min}} اقلب القرص المسطّح أو الخبز المُسطَّح. واخبزه المدة نفسها من الوجه الآخر. وبحسب ما يعجبك، يمكنك الانتظار قليلًا لتسمح للخبز بأن يصير محروقًا قليلًا. وكلما طالت مدة خبز خبزك، تبخّر منه الحموضة أكثر. وإن كان عجينك أكثر حامضية قليلًا، فيمكنك استخدام هذه الحيلة لموازنة الحموضة. وهذا يعتمد فعلًا على النكهة التي تسعى إليها.

عند صنع الخبز المُسطَّح أنصحك بلفّ الأقراص المُسطَّحة المخبوزة في منشفة مطبخ. فهذه الطريقة تجعل جزءًا أكبر من الرطوبة المتبخرة يبقى داخل خبزك، مما يضمن بقاء أقراصك المُسطَّحة طيبة ومنتفخة لمدة أطول بعد الخَبز. وتُستخدم استراتيجية مشابهة عند صنع تورتيلا الذرة.

يمكنك تخزين الأقراص المُسطَّحة أو الخبز المُسطَّح المخبوز بأمان في ثلاجتك لأسابيع. وعند التخزين احرص على وضعها في كيس بلاستيكي محكم الإغلاق كي لا تجف. وإن جفت، فانثر عليها بعض الماء وحمّصها. ستكون شبه جيدة كيوم خبزها الطازج.

احتفظ بقليل من عجينك غير المخبوز. فيمكنك استخدامه لصنع الدفعة التالية من الخبز أو الأقراص المسطّحة في اليوم التالي. وإن أردت الخبز بعد أيام قليلة، فأضف قليلًا من الماء والدقيق واحفظ هذا الخليط في ثلاجتك ما دمت تشاء.\footnote{ستبقى البادئة جيدة لأشهر. وإن توقعت أن تتركها مدة أطول، ففكر في تجفيف قليل من بادئة الخبز الحامض لديك.}

\subsection{وصفة الخبز المُسطَّح البسيطة}\label{subsec:flat-bread-recipe}

باتباع الخطوات المبينة في هذا القسم، ستتعرّف على خبز متعدد الاستخدامات يصلح لاستعمالات طبخية لا حصر لها. سواء كنت تغمسونه في صلصة مالحة، أو تلفّون فيه حشوة منكّهة، أو تتناولونه ببساطة مع رشّة زيت زيتون، فسيسعدك هذا الخبز المُسطَّح بلا شك.

% flat-bread-ar.tex — Arabic edition of upstream recipes/flat-bread.tex
% (input by ch. 8 at subsec:flat-bread-recipe, وصفة الخبز المُسطَّح البسيطة)
%
% Original: The Sourdough Framework by Hendrik Kleinwächter,
% the-sourdough-framework.com, CC BY-SA 4.0. The Arabic edition inherits
% CC BY-SA 4.0.
%
% Phase 28R R2 unit 10 reconstruction. Text recovered verbatim from the
% R0 recovery reference (recovery-r0/text/08-bread-types.txt, mined from
% the Phase 23 website build of the accepted Arabic sources) with the
% accepted edition book-ar-fixed.pdf (b78c2a28b574288801d8000a7249af3a)
% as the fidelity authority; no text retranslated from the English.
%
% Ingredients tabular keeps the upstream column spec {r@{}rl@{}}
% verbatim (NOT mirrored): the accepted render of p. 68 shows the
% ingredient column flush right at x520.9 with the amounts leftmost,
% exactly the upstream visual order. Every non-empty cell is wrapped
% {\textdir TRT ...} per the unit 9 cell-direction binding. Row 1's
% percent cell carries the upstream leading ~ (upstream writes
% "&~(" for the 100% row): inside the TRT cell it renders as a
% trailing visual space, which the accepted render shows (row 1's
% closing paren ends x281.3 vs rows 2-4 at x285.0); rows 2-4 have no ~.
% Documented residual (paren-font): the accepted edition renders these
% ASCII parens in Amiri; the lua fontsplit remap re-fonts ASCII
% punctuation to Pagella unconditionally, so our build renders them
% Pagella. Unfixable at the content layer; logged for the R3
% harmonization agenda.
%
% The الخطوات list is a description list per the established
% convention (starter-ar.tex): \item[label:] with the sty's
% style=nextline, leftmargin=0cm settings; the first item's label
% carries no colon, matching the upstream "Prepare the dough".
% \emph{ضبط زمن الطهي.} renders Amiri-Italic in the accepted render
% (verified p. 68), matching the upstream \emph{Adjusting cook time.}

\subsubsection*{المقادير}

\begin{tabular}{r@{}rl}
{\textdir TRT \LR{\qty{400}{\gram}}} & {\textdir TRT ~(\LR{\qty{100}{\percent}})} & {\textdir TRT دقيق (قمح، جاودار، ذرة، أيّ ما يكون متوفرًا لديك)} \\
{\textdir TRT \LR{\qty{320}{\gram}}} & {\textdir TRT (\LR{\qty{80}{\percent}})} & {\textdir TRT ماء، ويُستحسن أن يكون بحرارة الغرفة} \\
{\textdir TRT \LR{\qty{80}{\gram}}} & {\textdir TRT (\LR{\qty{20}{\percent}})} & {\textdir TRT بادئة خبز حامض نشيطة} \\
{\textdir TRT \LR{\qty{8}{\gram}}} & {\textdir TRT (\LR{\qty{2}{\percent}})} & {\textdir TRT ملح} \\
\end{tabular}

\subsubsection*{الخطوات}

\begin{description}
\item[اعجن العجين] في وعاء خلط كبير، امزج الدقيق والماء. اخلط حتى تحصل على عجين غير منتظم القوام لا بقع جافة فيه. أضف بادئة الخبز الحامض والملح إلى الخليط. وامزجها جيدًا حتى تحصل على عجين أملس ومتجانس.
\item[التخمير:] غطِّ الوعاء بغطاء أو بغلاف بلاستيكي. واترك العجين يرتاح ويتخمر حتى يزداد \LR{\qty{50}{\percent}} على الأقل في الحجم. وبحسب الحرارة ونشاط بادئتك، قد يستغرق ذلك ما بين \LR{\qty{4}{\hour}} و\LR{\qty{24}{\hour}}.
\item[تحضير الطهي:] حين ينتفخ العجين، سخّن مقلاة على حرارة متوسطة. وادهن المقلاة بزيت خفيف، مع مسح أي زيت زائد بمنشفة ورقية.
\item[التشكيل والطهي:] بمغرفة أو بيدك، اسحب جزءًا من العجين وضعه على المقلاة الساخن، ونثره برفق كما تفعل مع قرص مسطّح. غطِّ المقلاة بغطائها. فهو يحبس البخار ويضمن طهيًا متساويًا من الأعلى، ما يجعل قلبها لاحقًا أسهل. بعد نحو \LR{\qty{5}{\minute}}، أو حين يكتسب وجه الخبز المُسطَّح قشرة بنية ذهبية، اقلبه بحذر بسباتولا. \emph{ضبط زمن الطهي.} إذا بدا الخبز المُسطَّح داكنًا أكثر من اللازم، فتذكّر أن تقلل زمن الطهي قليلًا في القرص التالي. وبالعكس، إن كان شاحبًا، فدعه يطهى مدة أطول قليلًا قبل قلبه. اخبز الوجه المقلوب \LR{\qty{5}{\minute}} إضافية أو حتى يصير هو الآخر بنيًا ذهبيًا.
\item[التخزين:] بمجرد نضج الخبز، أخرجه من المقلاة وضعه على منشفة مطبخ. ولفّ الأقراص في المنشفة يساعد على حفظ طراوتها ويمنعها من أن تصير قاسية أكثر من اللازم. وكرر عملية الطهي لما تبقى من العجين.
\item[اقتراح للتقديم:] تناول أقراص الخبز الحامض المُسطَّحة ساخنة، مع الصلصات أو المربيات المفضلة لديك، أو بوصفها طبقًا جانبيًا لأي وجبة.
\end{description}


\section{خبز القالب}\label{sec:loaf-pan-bread}

يُصنع خبز القالب بمساعدة قالب خبز \textenglish{(loaf pan)} أو قالب صفيح \textenglish{(loaf tin)} خاص. وتوفر حواف القالب دعمًا إضافيًا للعجين كي ينتفخ. وصنع الخبز باستخدام قالب الخبز يتطلب فرنًا.

\begin{figure}[!htb]
\centering
\includegraphics[width=\textwidth]{loaf-pan-free-standing}
\caption[خبز قائم بذاته وخبز القالب]{خبز قائم بذاته وخبز قالب من القمح. وكلاهما تلقّى شقًّا صغيرًا قبل الخَبز يساعد على التحكم في طريقة انفتاحه.}\label{fig:free-standing-loaf-pan}
\end{figure}

بعد عجن عجينك، يمكنك وضعه مباشرة في قالب الخبز. وسيصير هذا أبسط العملية كلها إذ يمكنك تجاوز خطوات مثل تشكيل العجين.

لصنع خبز قالب رائع بجهد يسير:
%
\begin{enumerate}
    \item اخلط مقادير عجينك (الدقيق الخالي من الغلوتين يعمل أيضًا)
    \item ضعه في قالب الخبز
    \item انتظر حتى يتضاعف حجم عجينك تقريبًا
    \item اخبزه في فرن غير مُسخَّن مسبقًا لنحو \LR{\textenglish{30 min\textendash 50 min}}
\end{enumerate}

معرفة زمن الخَبز الدقيق أحيانًا شأن صعب قليلًا إذ قد يكون خارج خبزك قد نضج بينما الداخل ما يزال نيئًا. والطريقة الأفضل هي استخدام مقياس حرارة وقياس الحرارة اللُّبّية. فعند نحو \LR{\qty{92}{\degreeCelsius}} (\LR{\qty{197}{\degF}}) يكون عجينك قد جهز. وأنا أخبز خبز القالب عادةً عند نحو \LR{\qty{200}{\degreeCelsius}} (\LR{\qty{390}{\degF}})، وهي حرارة تقل قليلًا عن خبزي الرغيف القائم بذاته الذي أخبزه عند \LR{\qty{230}{\degreeCelsius}} (\LR{\qty{445}{\degF}}). وسبب ذلك أن العجين يحتاج إلى وقت حتى ينضج كما ينبغي داخل قالب الخبز. فالحواف لا تسخن بذلك السرعة. فينضج الجزء العلوي من العجين نضجًا جيدًا بينما الداخل لم ينضج بعد. وعند الخَبز احرص على استخدام بخار، أو ضع ببساطة قالب خبز آخر بمقاس قالبك نفسه فوق قالبك. وبهذا تحاكي القدر الهولندي. وستبقى رطوبة العجين المتبخرة داخله.

وحيلة جيدة لصنع خبز قالب ممتاز هي صنع عجين لزج جدًا. يمكنك اختيار ترطيب بقيمة \LR{\textenglish{90 \%\textendash 100 \%}}، يشبه تقريبًا بادئة الخبز الحامض الافتراضية. ومثل الخبز المُسطَّح، تساعد الرطوبة العالية على صنع لُبّ أكثر هوائية وانتفاخًا. وسيكون الخبز أيضًا أكثر طعمًا مطاطيًا في القرص. وهذا النوع من الخبز المصنوع من الجاودار هو نمط عائلتي المفضل. فنكهة الجاودار القوية مع القوام اللزج يصنعان حقًا خبز ساندويتش ممتازًا.

ولتحسين البنية يمكنك أيضًا التفكير في استخدام نحو \LR{\qty{50}{\percent}} من دقيق القمح في خلطتك. فشبكة الغلوتين ستتكون أثناء تخمر عجينك وستسمح بحبس غاز أكبر في العجين.

ومشكلة شائعة ستواجهها عند صنع خبز القالب هي تصاق العجين بالقالب. استخدم كمية كبيرة من الزيت لتدهين قالبك. ويمكن لبخاخ زيت نباتي غير لاصق أن يعمل العجائب. لا تنظف قوالب الخبز بالصابون. استخدم منشفة مطبخ فقط لتنظيفها. ومع كل خَبز تتكون طبقة أفضل، فيصير قالبك أكثر مقاومة للالتصاق.

والرائع في هذا النوع من الخبز أنه يعمل مع كل دقيق. والزمن الكلي لمعالجة العجين ربما لا يتجاوز \num{5} دقائق، مما يجعل دمجه في روتينك اليومي سهلًا جدًا. كما أن قوالب الخبز تستغل المساحة في فرنك بكفاءة عالية. وباستخدام القوالب يمكنني خبز \num{5} أرغفة في الوقت نفسه بسهولة في فرني المنزلي. أما خبز الأرغفة القائمة بذاتها فأحتاج له عادةً إلى جلسات خَبز متعددة.

\section{الخبز القائم بذاته}

يُخبز الرغيف القائم بذاته كليًا دون أوعية خبز داعمة في فرنك. ولصنع رغيف قائم بذاته تلزم خطوات وأدوات أكثر.

\begin{figure}[!htb]
\centering
\includegraphics[width=0.85\textwidth]{free-standing-loaf}
\caption[خبز حامض قائم بذاته]{خبز حامض قائم بذاته. لاحظ أن الشق المعروف بـ\emph{الأذن} \textenglish{(ear)} والانتفاخ في الفرن \textenglish{(oven spring)} يميزان بوضوح هذا النوع من الخبز عن الخبز المُسطَّح وخبز القالب.}
\end{figure}

عند استخدام القمح، احرص على عجن عجينك بما يكفي لتكوين شبكة غلوتين. واترك للعجين أن يبلغ زيادة معينة في الحجم أثناء التخمير. وبعد ذلك قسّم العجين وشكّله مسبقًا \textenglish{(pre-shape)} بالشكل الظاهري الذي تريده. فكل شكل يتطلب تقنية مختلفة. وأحيانًا يكون بلوغ الشكل الصحيح صعبًا. فصنع خبز الباجيت \textenglish{(baguette)} مثلًا، يتطلب أداء خطوات أكثر. وإتقان هذه التقنية يستهلك عدة محاولات.

وبعد تشكيل العجين، يُترك للتخمير النهائي مجددًا لمدة معينة. وحين يجهز العجين، تُستخدم أداة حادة مثل شفرة الحلاقة لإحداث شق في العجين. فهذا يساعد على التحكم في طريقة انفتاح العجين أثناء عملية الخَبز.

كل هذه الخطوات تتطلب تدريبًا. وينبغي أداء كل منها على وجه الإتقان، دون أخطاء. لكن بعد الخَبز سيكافئك خبز جميل بطعم وملمس رائعين.

وهناك وصفة وشرح مخصصان لهذا النوع من الخبز في فصل الخبز الحامض القمحي.
